\section{Probleme}

\subsection{Problema Carry Bit (IIOT 2023-2024 runda 3)}
\label{problem:carry-bit}

\href{https://kilonova.ro/problems/2165}{enunț}
$\bullet$
\hyperref[code:carry-bit]{sursă}

\subsubsection*{Observații preliminare}

Să considerăm două numere binare de lungime $l$, $p[0 \dots l-1]$ și $q[0 \dots l-1]$ (unde $p[0]$ și $q[0]$ sînt biții cei mai semnificativi). Cînd există depășire la adunare?

\begin{itemize}
  \item Dacă $p[0] = q[0] = 1$, atunci în mod evident există depășire.

  \item Dacă $p[0] = q[0] = 0$, atunci nu există depășire. Ambele numere sînt mai mici decît $2^{l-1}$, deci suma lor este mai mică decît $2^l$.

  \item Dacă $p[0] \neq q[0]$, atunci există depășire dacă și numai dacă există depășire la adunarea $p[1 \dots n-1] + q[1 \dots n-1]$.
\end{itemize}

Atunci, pentru problema dată, pentru interogarea $\langle x, y, l \rangle$ trebuie să aflăm prima poziție la care șirurile $a[x \dots]$ și $b[y \dots]$ coincid. Există depășire dacă și numai dacă (1) poziția este la distanță cel mult $l$ de pozițiile $x$, respectiv $y$, și (2) valorile care coincid sînt egale cu 1.

În practică, este mai simplu să negăm unul dintre șiruri, să spunem $b$, deoarece problema găsirii primei \textbf{diferențe} între două subșiruri este mai abordabilă. Notăm acest șir cu $\bar{b}$.

\subsubsection*{Soluția cu hashing}

Menționez că există și o soluție mai simplă, bazată pe hashing (\href{https://kilonova.ro/submissions/176125}{exemplu}). Interpretăm șirurile ca pe numere într-o bază arbitrară, dar fixă. Nu baza 2, deoarece este foarte probabil că există teste contra ei. Pentru fiecare poziție $i$ calculăm valoarea hash a prefixului $a[0 \dots i]$. Atunci prin diferență de prefixe putem afla valoarea hash a oricărui subșir al lui $a$, și similar pentru $\bar{b}$.

Dată fiind o interogare, căutăm binar prima lungime $w$ la care valorile hash pentru $a[x \dots x + w - 1]$ și $\bar{b}[y \dots y + w - 1]$ diferă. Apoi considerăm simplu cazurile descrise mai sus.

\subsubsection*{Soluția cu șiruri de sufixe}

Dacă concatenăm cele două șiruri (separate de un delimitator ca \texttt{\textbackslash0}), atunci putem reformula interogările astfel: date fiind două poziții $x$ și $y$, care este lungimea maximă a două subșiruri egale care încep de la acele poziții? Această problemă este aparent clasică și se numește \textit{longest common extension} (\textit{LCE}).

Dacă nu ne vine ideea cu hashing, nu strică să avem și arma universală a șirurilor de sufixe. Construim cele trei informații necesare:

\begin{enumerate}
  \item Șirul de sufixe.
  \item Vectorul $LCP$.
  \item O structură pentru RMQ. Eu am folosit un AINT.
\end{enumerate}

Acum, date fiind pozițiile de start $x$ și $y$, aflăm cu RMQ lungimea $w$ a celui mai lung prefix comun al sufixelor care încep la pozițiile $x$ și $y$. Iau naștere aceleași trei cazuri ca mai sus:

\begin{enumerate}
  \item Dacă $w \geq l$, atunci $a$ și $\bar{b}$ coincid pe toată lungimea $l$. Echivalent, $a$ și $b$ diferă pe toată lungimea $l$, deci suma lor este exact $2^{l-1}$ și nu există depășire.

  \item Altfel, dacă $a[x + w] = b[y + w] = 0$, atunci nu există depășire.

  \item Altfel $a[x + w] = b[y + w] = 1$ și există depășire.
\end{enumerate}


\subsection{Problema Periodic Words (IIOT 2023-2024 runda 1)}
\label{problem:periodic-words}

\href{https://kilonova.ro/problems/1940}{enunț}
$\bullet$
\hyperref[code:periodic-words]{surse}

Problema pornește de la următoarea observație. Pentru o lungime de perioadă aleasă, putem reduce testul de periodicitate la o singură comparație de sufixe. Sper că v-ați amintit acest truc de la problema \hyperref[problem:bart]{Bart}.

De exemplu, pentru a afla dacă $S[20 \dots 79]$, de lungime 60, are perioadă 5, vrem să aflăm dacă subșirurile $S[20 \dots 24], S[25 \dots 29], \dots S[74 \dots 79]$ sînt egale. Dar este suficient să testăm dacă $S[20 \dots 74] = S[25 \dots 79]$.

Astfel, ambele soluții de mai jos calculează, pentru fiecare interogare $[l, r]$ și fiecare divizor $d$ al lui $r - l + 1$, dacă $S[l \dots r - d] = S[l + d \dots r]$. Ca optimizare, este suficient să testăm doar perioadele care corespund unui număr de repetiții prim. De exemplu, dacă un subșir de lungime 60 poate fi scris ca 20 de repetiții cu perioadă 3, el poate fi scris și ca 2 repetiții cu perioadă 30 sau 5 repetiții cu perioadă 12.

\subsubsection*{Soluția cu hashing}

O primă soluție calculează hash-urile prefixelor într-o bază $B$ și modulo o valoare $M$. Din acestea putem calcula hash-ul oricărui subșir.

Întrucît există circa $5 \cdot 10^9$ subșiruri, paradoxul zilei de naștere ne spune că trebuie să operăm modulo o valoare cel puțin $25 \cdot 10^{18}$ și preferabil chiar mai mare. Deci, la prima vedere, trebuie să alegem două module pentru a elimina coliziunile. De fapt, situația este mai simplă. Noi nu comparăm decît subșiruri \textbf{de aceeași lungime}, deci avem cel mult $10^6$ subșiruri în aceeași clasă. Coliziunile între subșiruri de lungimi diferite nu ne deranjează. Dacă operăm modulo $2^{64} \approx 1,8 \cdot 10^{19}$, probabilitatea unei coliziuni este infimă.

\subsubsection*{Soluția cu șiruri de sufixe}

Calculăm șirul de sufixe al lui $S$, precum și vectorul $LCP$ și informații de RMQ peste acesta. Acum, date fiind interogarea $[l, r]$ și divizorul $d$, verificăm dacă sufixele care încep la pozițiile $l$ și $l + d$ au un prefix comun de lungime cel puțin $r - l + 1 - d$.


\subsection{Problema Phone Book (Finala IIOT 2022-2023)}
\label{problem:phone-book}

\href{https://kilonova.ro/problems/457}{enunț}
$\bullet$
\hyperref[code:phone-book]{sursă}

Această problemă este relativ „servită” pentru șiruri de sufixe.

Cînd luăm în calcul un subșir al cuvîntului $i$ și ne întrebăm dacă este specific acelui cuvînt, trebuie să-l putem căuta rapid în toate celelalte cuvinte. De aceea, să construim șirul de sufixe al concatenării tuturor cuvintelor: $S = s_1 \texttt{\textbackslash0} \; s_2 \texttt{\textbackslash0} \; \dots \texttt{\textbackslash0} \; s_n \; \texttt{\textbackslash0}$.

Acum, să considerăm o poziție de start în interiorul cuvîntului $i$. Ne întrebăm: ca să obținem un subșir specific cuvîntului $i$, cîte caractere ar trebui să selectăm? Să spunem că el are poziția $k$ în vectorul de sufixe. \textbf{În principiu}, vrem să diferențiem subșirul de precedentul și de următorul din vectorul de sufixe. Deci trebuie să selectăm $1 + \max(LCP[k], LCP[k + 1])$ caractere.

Situația se complică un pic deoarece nu ne deranjează dacă acel sufix mai apare altundeva \textbf{în cuvîntul $i$}. Deci trebuie să căutăm înainte și înapoi prin vectorul $LCP$ și să găsim cele mai apropiate sufixe care nu provin din cuvîntul $i$. Nu putem face naiv (liniar) această căutare, deoarece ea duce la timp pătratic dacă există puține cuvinte lungi. Dar putem procesa grupe de sufixe consecutive care provin din același cuvînt.

Concret, pentru fiecare poziție din șirul de sufixe mai reținem două informații: (1) din ce cuvînt provine acel sufix și (2) care este lungimea sufixului (pînă la următorul delimitator \texttt{\textbackslash0}). Construim vectorul $LCP$ și informații de RMQ peste $LCP$.

Acum căutăm grupe maximale de sufixe consecutive (lexicografic) care provin din același cuvînt. Dacă sufixele de la poziția $a$ la poziția $b$ toate provin din același cuvînt $i$, atunci pentru toate acestea sufixele relevante, precedent și următor, se află la pozițiile $a - 1$ și $b + 1$. De aceea, pentru fiecare poziție $a \leq k \leq b$, procedăm astfel:

\begin{enumerate}
  \item Calculăm $r_1 = RMQ(a - 1, k)$.

  \item Calculăm $r_2 = RMQ(k, b + 1)$.

  \item Calculăm $r = 1 + \max(r_1, r_2)$. Acesta este numărul de caractere necesare pentru a diferenția sufixul al $k$-lea, care provine din cuvîntul $i$, de orice alte subșiruri din orice alte cuvinte în afară de $i$.

  \item Dacă sufixul curent are cel puțin $r$ caractere (pînă la sfîrșitul cuvîntului), atunci minimizăm răspunsul pentru cuvîntul $i$ cu valoarea $r$.

  \item Altfel, sufixul curent nu are destule caractere pentru a-l putea diferenția de celelalte cuvinte.
\end{enumerate}

\subsubsection*{Optimizare}

Putem omite cu totul structura de date pentru RMQ. Să observăm, de fapt, că toate interogările de minim sînt grupate. Pentru fiecare poziție $a \leq k \leq b$ ne interesează minimul pe intervalele $[a - 1, k]$ și $[k, b + 1]$. Putem calcula aceste minime prin două treceri, una înapoi și una înainte, prin intervalul $[a, b]$.
